\section{研究背景与意义}

随着移动互联网、工业互联网与智能终端的普及，图像数据已经成为信息系统中最主要的数据形态之一。如何从海量图像中快速、准确地提取语义信息，是智慧校园、智慧交通和智能制造等应用场景的关键问题。传统图像识别方法通常依赖人工设计特征，再结合分类器完成识别任务，这类方法在小规模数据与规则场景中具有一定效果，但面对复杂背景、尺度变化和遮挡噪声时往往难以保持稳定性能。

深度学习模型尤其是卷积神经网络在 ImageNet 竞赛中的突破性表现，推动图像识别技术进入数据驱动的新阶段。\supercite{krizhevsky2012imagenet,he2016deep} 相较于传统方法，深度神经网络能够通过端到端训练自动学习分层特征表示，使得模型在特征表达能力与泛化能力方面都取得显著提升。对于本科毕业设计而言，围绕一个清晰的视觉任务展开模型设计、实验对比和结果分析，既符合学校对于论文结构和实验论证的要求，也有利于训练学生的问题拆解与工程实现能力。

\section{国内外研究现状}

\subsection{传统图像识别方法}

传统方法通常遵循“特征提取 + 分类器”的研究范式。尺度不变特征变换、方向梯度直方图等方法通过人工经验描述图像局部纹理和边缘信息，再将提取后的特征输入支持向量机、随机森林等分类器中完成识别。\supercite{lowe2004sift,dalal2005hog} 该类方法的优势在于可解释性较强、计算开销相对可控，但局部描述子难以适应复杂环境下的语义变化，也难以覆盖高层抽象特征。

\subsection{深度学习图像识别方法}

AlexNet、VGGNet 和 ResNet 等模型的提出标志着深度学习在图像识别任务中的持续突破。\supercite{krizhevsky2012imagenet,simonyan2015vgg,he2016deep} 研究者进一步通过注意力机制、多尺度特征融合和高效卷积设计改善模型性能。\supercite{hu2018senet,tan2019efficientnet} 从整体趋势来看，图像识别模型的发展逐渐从“单纯追求更深更大”转向“精度、效率与可部署性并重”，这也是当前本科毕业设计中较为适合切入的研究方向。

\section{研究内容与论文结构}

本文以轻量化图像识别模型为研究对象，主要工作包括：

\begin{enumerate}
  \item 梳理图像识别领域的典型研究脉络，总结卷积神经网络与注意力机制的关键思想。
  \item 设计一类融合通道注意力和多尺度特征聚合的改进模型，用于提升复杂场景下的特征表达能力。
  \item 构建实验评测流程，从准确率、参数量与推理效率等角度对模型进行验证与分析。
  \item 总结研究结论，讨论工作局限，并提出后续改进方向。
\end{enumerate}

全文结构安排如下：第一章介绍研究背景、研究意义与论文结构；第二章阐述相关理论与关键技术；第三章给出模型设计思路与实现方法；第四章展示实验设计、结果比较与误差分析；第五章总结全文并对未来工作进行展望。
